\DocumentMetadata{
  pdfversion=2.0,
  pdfstandard=ua-2,
  lang=en-US,
  tagging=on,
  tagging-setup={math/setup=mathml-SE,tabsorder=structure}
}
\documentclass[11pt,letterpaper]{article}
\usepackage[margin=0.68in,includefoot]{geometry}
\usepackage{newcomputermodern}
\usepackage{amsmath,amscd,mathtools}
\usepackage{tikz,adjustbox}
\providecommand{\square}{\mdlgwhtsquare}
\usepackage[hidelinks]{hyperref}
\hypersetup{
  pdftitle={Numerical Analysis PhD Exam, May 2006},
  pdfauthor={University of Florida Department of Mathematics},
  pdfsubject={Graduate mathematics examination},
  pdfdisplaydoctitle=true,
  bookmarksopen=true
}
\setcounter{secnumdepth}{0}
\setlength{\parindent}{0pt}
\setlength{\parskip}{0.28em}
\setlength{\emergencystretch}{3em}
\setlength{\leftmargini}{2.15em}
\setlength{\leftmarginii}{2.65em}
\setlength{\leftmarginiii}{3em}
\pagestyle{empty}
\raggedbottom
\allowdisplaybreaks
\NewTaggingSocketPlug{block/list/label}{examproblem}{%
  \tagstructbegin{tag=Lbl}\tagstructend%
  \tagstructbegin{tag=\LBody}%
  \tagstructbegin{tag=\ProblemHeadingTag,title-o={\ProblemHeadingText}}%
  \tagmcbegin{tag=\ProblemHeadingTag,actualtext={\ProblemHeadingText}}%
  #2%
  \tagmcend%
  \tagstructend%
}
\newenvironment{examproblems}[2]{%
  \def\ProblemHeadingTag{#2}%
  \AssignTaggingSocketPlug{block/list/label}{examproblem}%
  \begin{enumerate}%
  \setcounter{enumi}{#1}%
  \renewcommand{\labelenumi}{\textbf{\theenumi.}}%
  \setlength{\topsep}{0.35em}%
  \setlength{\partopsep}{0pt}%
  \setlength{\itemsep}{0.48em}%
  \setlength{\parsep}{0.18em}%
}{\end{enumerate}}
\newenvironment{examsubparts}{%
  \AssignTaggingSocketPlug{block/list/label}{default}%
  \begin{enumerate}%
  \setlength{\topsep}{0.18em}%
  \setlength{\partopsep}{0pt}%
  \setlength{\itemsep}{0.16em}%
  \setlength{\parsep}{0.1em}%
}{\end{enumerate}}
\newenvironment{examinstructions}{%
  \par\begingroup\itshape
}{%
  \par\endgroup\vspace{0.18em}
}
\makeatletter
\renewcommand\subsection{\@startsection{subsection}{2}{0pt}%
  {-1.25ex plus -.25ex minus -.1ex}%
  {0.55ex plus .1ex}%
  {\normalfont\large\bfseries}}
\renewcommand{\@maketitle}{%
  \newpage\null\vskip -1em%
  {\footnotesize\bfseries
    \href{https://www.ufl.edu/}{University of Florida}, \href{https://math.ufl.edu/}{Department of Mathematics}\par}%
  \vskip -0.5em%
  \tagstructbegin{tag=H1,title={Numerical Analysis PhD Exam, May 2006}}%
  \tagpdfparaOff%
  \tagmcbegin{tag=H1}%
    {\LARGE\bfseries\href{https://gma.math.ufl.edu/exams/numerical-analysis-phd/}{Numerical Analysis PhD Exam}, May 2006\par}%
  \tagmcend%
  \tagpdfparaOn%
  \tagstructend%
  \vskip 0.25em%
  \tagmcbegin{artifact}%
    \hrule height 1pt%
  \tagmcend%
  \vskip 1em%
}
\makeatother
\title{Numerical Analysis PhD Exam, May 2006}
\author{}
\date{}
\begin{document}
\maketitle
\thispagestyle{empty}
\begin{examinstructions}
Do any eight of the ten problems below.
\end{examinstructions}
\begin{examproblems}{0}{H2}
\def\ProblemHeadingText{Problem 1.}
\item
Let \(x_0,\ldots,x_n\) be distinct real points. Let \(\{l_j(x)\}_{j=0,\ldots,n}\) be the Lagrange basis functions for these points. Prove that
\[
\sum_{j=0}^n l_j(x)=1
\]
 for all~\(x\).
\def\ProblemHeadingText{Problem 2.}
\item
Suppose you want to solve numerically the ode
\[
y'=f(x,y),\qquad y(0)=y_0.
\]
\begin{examsubparts}
\item[\textnormal{(a)}]
By applying the midpoint rule to the integral in the equation
\[
y(x_{n+1})=y(x_{n-1})+\int_{x_{n-1}}^{x_{n+1}}f(s,y(s))\,ds,
\]
 derive the midpoint method for odes.
\item[\textnormal{(b)}]
Find the order of the local truncation error for this method.
\item[\textnormal{(c)}]
Determine if this method is strongly stable, weakly stable or unstable.
\end{examsubparts}
\def\ProblemHeadingText{Problem 3.}
\item
Given \(x_0,x_1\) and smooth \(f(x)\), give a divided difference type formula for the cubic polynomial which satisfies:
\[
\begin{aligned}
p(x_0)&=f(x_0),\\
p'(x_0)&=f'(x_0),\\
p''(x_0)&=f''(x_0),\\
p(x_1)&=f(x_1).
\end{aligned}
\]
 Also give a divided difference type error formula for \(f(x)-p(x)\).
\def\ProblemHeadingText{Problem 4.}
\item
This problem has two parts.
\begin{examsubparts}
\item[\textnormal{(a)}]
Find the linear least squares approximation to \(f(x)=x^2\) on the interval \([0,2]\) by optimizing the least squares error over~\(a\) and~\(b\) in \(r^*(x)=ax+b\).
\item[\textnormal{(b)}]
Find the first two orthonormal polynomials \(\phi_0(x),\phi_1(x)\) on \([0,2]\) (weight function \(w(x)=1\)). Use these to find the linear least squares approx to \(f(x)=x^2\), and check this result agrees with the previous problem.
\end{examsubparts}
\def\ProblemHeadingText{Problem 5.}
\item
Consider the root-finding iteration defined by
\[
x_{n+1}=x_n-f(x_n)\left[\frac{f(x_n)}{f(x_n+f(x_n))-f(x_n)}\right].
\]
 This is called Steffensen’s method. Show that the method converges quadratically when applied to \(f(x)=x^2-a\). (You may assume \(x_0\) is close enough to the root.)
\def\ProblemHeadingText{Problem 6.}
\item
This problem has three parts.
\begin{examsubparts}
\item[\textnormal{(a)}]
Evaluate the~\(p\)-norm of a diagonal matrix, \(1\leq p\leq\infty\).
\item[\textnormal{(b)}]
Evaluate the~\(p\)-norm of a rank one matrix \(\mathbf u\mathbf v^*\), \(1\leq p\leq\infty\), \(\mathbf u\in\mathbb C^m\) and \(\mathbf v\in\mathbb C^n\).
\item[\textnormal{(c)}]
Suppose \(\mathbf A\in\mathbb C^{m\times n}\) can be expressed
\[
\mathbf A=\sum_{k=1}^r\sigma_k\mathbf u_k\mathbf v_k^*,
\]
 where the vectors \(\{\mathbf u_k:1\leq k\leq r\}\) and \(\{\mathbf v_k:1\leq k\leq r\}\) are orthonormal. What is \(\|\mathbf A\|_2\)?
\end{examsubparts}
\def\ProblemHeadingText{Problem 7.}
\item
This problem has two parts.
\begin{examsubparts}
\item[\textnormal{(a)}]
Given \(\mathbf A\in\mathbb C^{m\times n}\), show that \(\mathbf A^*\mathbf A\) is positive semidefinite and \(\mathbf A^*\mathbf A\) is positive definite if and only if the columns of \(\mathbf A\) are linearly independent.
\item[\textnormal{(b)}]
If \(\mathbf A\) is a Hermitian, positive semidefinite matrix, show that the diagonal contains the largest in magnitude element of \(\mathbf A\). Hint: show that \(\lvert a_{ij}\rvert\leq\max\{\lvert a_{ii}\rvert,\lvert a_{jj}\rvert\}\).
\end{examsubparts}
\def\ProblemHeadingText{Problem 8.}
\item
Let \(\mathbf A\in\mathbb C^{n\times n}\) be nonsingular. Show that \(\mathbf A\) has an LU factorization if and only if for each~\(k\) with \(1\leq k\leq n\), the upper-left \(k\times k\) block \(\mathbf A_{1:k,1:k}\) is nonsingular. Prove that this LU factorization is unique.
\def\ProblemHeadingText{Problem 9.}
\item
Consider the Arnoldi iteration on the Krylov spaces \(\mathcal K_k=\operatorname{span}\{\mathbf b,\mathbf A\mathbf b,\mathbf A^2\mathbf b,\ldots,\mathbf A^{k-1}\mathbf b\}\) for some \(\mathbf b\in\mathbb C^m\) and \(\mathbf A\in\mathbb C^{m\times m}\) (given at side). Suppose the algorithm proceeds without breakdown until for some \(n<m\), it encounters \(h_{n+1,n}=0\).

\par
Algorithm 1 (Arnoldi iteration)

\par
(a) Set \(\mathbf q_1=\mathbf b/\|\mathbf b\|_2\).

\par
(b) For \(n=1,2,3,\ldots\) do:

\par
i. Set \(\mathbf v=\mathbf A\mathbf q_n\).

\par
ii. For \(j=1,2,\ldots,n\) do:

\par
A. Set \(h_{jn}=\mathbf q_j^*\mathbf v\).

\par
B. Replace \(\mathbf v\) by \(\mathbf v-h_{jn}\mathbf q_j\).

\par
iii. Set \(h_{n+1,n}=\|\mathbf v\|_2\).

\par
iv. Set \(\mathbf q_{n+1}=\mathbf v/h_{n+1,n}\).
\begin{examsubparts}
\item[\textnormal{(a)}]
Show that \(\mathcal K_n\) is an invariant subspace of \(\mathbf A\), i.e., \(\mathbf A\mathcal K_n\subseteq\mathcal K_n\).
\item[\textnormal{(b)}]
Show that \(\mathcal K_n=\mathcal K_{n+1}=\mathcal K_{n+2}=\cdots\).
\item[\textnormal{(c)}]
Let \(\mathbf H_n\) be the \(n\times n\) Hessenberg matrix whose \(ij\)-th entry (\(i\leq j+1\)) is the number \(h_{ij}\) computed in the algorithm. Show that each eigenvalue of \(\mathbf H_n\) is an eigenvalue of \(\mathbf A\).
\end{examsubparts}
\def\ProblemHeadingText{Problem 10.}
\item
Consider an invertible linear system \(\mathbf A\mathbf x=\mathbf b\) and a perturbation \((\mathbf A+\delta\mathbf A)(\mathbf x+\delta\mathbf x)=\mathbf b\). Assuming \(\|\mathbf A^{-1}\|\|\delta\mathbf A\|<1\), derive the condition estimate
\[
\frac{\|\delta\mathbf x\|}{\|\mathbf x\|}\leq\left(\frac{\|\mathbf A\|\|\mathbf A^{-1}\|}{1-\|\mathbf A^{-1}\|\|\delta\mathbf A\|}\right)\frac{\|\delta\mathbf A\|}{\|\mathbf A\|}.
\]
 For given \(\mathbf A\) and \(\mathbf x\) and for the 2-norm, what choice of \(\delta\mathbf A\) yields the following equality:
\[
\|\mathbf A^{-1}\delta\mathbf A\mathbf x\|=\|\mathbf A^{-1}\|\|\delta\mathbf A\|\|\mathbf x\|?
\]
\end{examproblems}
\end{document}
